\chapter{Passive Radar Processing}

The passive-radar chain equalizes reference and surveillance channels, cancels direct signal and clutter, computes the cross-ambiguity function, detects with false-alarm control and associates detections across nodes. The output is a track with uncertainty, not only a coordinate. Direct-emitter detection must remain distinct from bistatic-reflection detection.

\section{Aircraft broadcast consistency experiment}
The reference aircraft signal is ADS-B, primarily 1090ES at 1090 MHz (\texttt{EXP-001}) for commercial aviation. The alternative 978 MHz UAT channel (\texttt{EXP-002}) is contextual and relevant mainly to certain general-aviation deployments in the United States. Both are assigned to the dedicated 960--1215 MHz aviation aperture or array rather than to the 470--860 MHz UHF Yagi. A diplexed receiver provides simultaneous independent service paths, each with its own preselector, limiter/LNA, coherent ADC and FPGA channel. ADS-B messages are transmitted automatically and repeatedly, with message timing dependent on the information type; they are not one continuous packet carrying every parameter at every instant. Expected fields include GNSS-derived position, altitude, velocity, track and aircraft identification.

Two radome nodes (\texttt{EXP-005}) separated by a proposed $100\,\mathrm{km}$ baseline (\texttt{EXP-003}) receive the aircraft transmission independently. This illustrative geometry is distinct from the three-node demonstrator target (\texttt{ARC-002}). For each node, the system records received power $P_i$, local angle of arrival $\mathrm{AOA}_i$, hardware timestamp $t_i$ and Doppler estimate $f_{D,i}$. The fusion layer computes:
\[
\Delta t=t_A-t_B,\qquad \Delta f_D=f_{D,A}-f_{D,B}.
\]
Together with surveyed node coordinates, aircraft ephemeris hypotheses and propagation uncertainty, these observables produce a consistency region for altitude, position and velocity. The result checks whether the broadcast parameters are compatible with the physical observation; it is not a blind guarantee that every transmitted parameter is truthful.

ADS-B supplies a cooperative identity and state report, while the two-node RF geometry supplies an independent consistency check based on received power, AOA, TDOA and Doppler.

The nominal 100 km baseline remains a screening geometry, not an approved deployment distance. With the proposed screening heights of 1000 m for a station and 10000 m for an aircraft (\texttt{EXP-010}), a standard effective-Earth radio-horizon calculation gives approximately 542 km of mutual geometric horizon. This only shows that 100 km is not excluded by curvature in that example. Terrain masking, antenna patterns, link margin, common visibility, GDOP, spectrum authorization and actual site coordinates must still close before field deployment.

\section{Independent UHF illuminators}
Known UHF television or other broadcast transmitters can provide an independent triangulation test alongside ADS-B. Their direct signals can be used to calibrate AOA, timing, received-power mapping and site geometry when transmitter coordinates and waveform characteristics are known. Their reflected signals can also support a separate bistatic experiment, but direct-emitter verification and target-reflection localization must not be mixed in the same estimator without an explicit propagation model.

The experiment should compare: (i) ADS-B at 1090 MHz as a cooperative direct emitter; (ii) known UHF broadcast signals, provisionally in the 470--860 MHz region (\texttt{EXP-004}), as direct reference illuminators; and (iii) bistatic echoes from a controlled or independently tracked target. Agreement across these channels tests the radome network's triangulation, calibration and covariance reporting rather than merely repeating the aircraft message.

\section{Independent experiment protocols}
The cooperative direct-emitter protocol (\texttt{EXP-006}) records decoded identity and state, received power, AOA, hardware time and Doppler. Its truth source is the decoded GNSS-derived state with independently surveyed receiver coordinates; its estimator tests RF-versus-message consistency, and its metrics are association rate and position, velocity, AOA, time and Doppler residuals with declared uncertainty.

The known-transmitter calibration protocol (\texttt{EXP-007}) records the direct UHF waveform, received power, AOA and calibrated channel delay. Its truth source is an independently surveyed transmitter position, antenna sector and waveform metadata; its estimator solves the expected direct path, and its metrics are AOA, delay and power-map residuals plus repeatability over time and temperature.

The bistatic-reflection protocol (\texttt{EXP-008}) uses simultaneous reference and surveillance channels, direct-path cancellation and a cross-ambiguity estimator. Its truth source is a controlled or independently tracked target trajectory plus surveyed transmitter and receiver coordinates; its metrics are probability of detection, false-alarm probability and delay, Doppler and localization errors with covariance consistency. No result from the two direct-signal protocols substitutes for this bistatic evidence.

\section{Bistatic ellipsoid and uncertainty-region reduction}
Let the illuminator position be $\mathbf{s}$, the surveyed position of receiver $i$ be $\mathbf{r}_i$, and a hypothesized target position be $\mathbf{x}$. The excess delay between the echo and the direct reference path satisfies
\[
c\,\Delta\tau_i=
\lVert\mathbf{s}-\mathbf{x}\rVert+
\lVert\mathbf{x}-\mathbf{r}_i\rVert-
\lVert\mathbf{s}-\mathbf{r}_i\rVert.
\]
For fixed $\Delta\tau_i$, compatible points form a bistatic iso-range ellipsoid whose foci are the transmitter and receiver. Per-face received power, corrected for the embedded pattern, losses and the illuminator's effective radiated power, supplies an angular/link likelihood; AOA restricts the solution to a direction or cone. Polarization observables add discrimination only at the same frequency and through calibrated coherent ports. The present VHF and UHF Yagis operate in different bands and therefore cannot be crossed to form a polarimetric measurement of the same echo.

The probabilistic intersection of delay ellipsoids, AOA, Doppler/FDOA, power and valid polarization yields a state region with covariance rather than an exact point. Associated reception of the same echo by multiple antennas and faces of one radome adds angular, power, delay and Doppler observations to the estimator. Reception of that echo by a second synchronized radome adds another ellipsoid and further constraints; additional integrated nodes add independent information and help reject ambiguities.

In weighted-least-squares reconciliation, or an equivalent probabilistic estimator, every additional valid observation adds a positive-semidefinite information matrix. The estimated covariance therefore cannot increase, and measurement uncertainty decreases strictly along directions for which the new observation supplies non-redundant information. This joint propagation of input quantities, covariances and multivariate output quantities follows the metrological treatment in GUM Supplement 2, JCGM 102:2011 \citep{jcgm1022011}. It is a mathematical consequence of the fusion model, not a claim of already measured performance. Prior knowledge of the illuminator position and waveform, reference channels, hardware timestamping, and the common time base---centrally referenced and continuously calibrated through GNSS, frequency reference, White Rabbit, holdover and delay measurement---allow the observations to be associated and reconciled. The demonstrator must quantify the magnitude of the reduction and compare predicted covariance with observed error and GDOP; geometry, SNR, multipath, calibration and error correlation determine the achieved gain, not the existence of the reduction principle.

\section{Jamming as an observable emission}
Electronic attacks against active radars may use noise, sweep, frequency agility, and coherent or deceptive waveforms. To RADOME, this energy does not disappear: it is an additional direct emission in the spectrum and increases passive detectability of the transmitter when it lies within the covered bands and dynamic range, especially at high power, bandwidth or duty cycle. NATO literature treats passive radar, passive emitter localization and coherent electronic attack against sensors---including PCL---as related but distinct problems \citep{nato2019passiveradar,nato2017coherentea}.

Protocol \texttt{EXP-011} treats jamming as a non-cooperative source. Multiple faces and radomes record power, spectrum/occupancy, AOA, timestamps and FDOA/Doppler; temporal and spectral association support estimation of the emission centroid and, through a sequence of states, its motion vector. Increasing radiated energy to obscure active radars therefore supplies the passive network with more direct energy to intercept and geolocate. Processing must preserve raw data before interference suppression, separate jammer detection/classification from the bistatic-echo chain, and use limiters, preselection and gain control to prevent saturation.

This advantage does not automatically identify the aircraft. Directional, low-probability-of-intercept, coherent, towed, expendable or stand-off jammers can create multiple sources, false angles and ambiguous association; the localized emitter may not be aboard the fighter. The controlled claim is that jamming increases the electromagnetic observability of its source to the passive network and creates additional observables. Tests must measure probability of intercept, AOA/TDOA/FDOA error, saturation, classification and track consistency for each signal class.




\section{Cellular towers as calibrated illuminators}
Cellular base stations add a particularly valuable set of known illuminators. Their surveyed coordinates, sector orientations and nominal carrier allocations allow the network to use multiple cellular frequencies and waveforms as direct references. The direct signal can calibrate coverage, AOA, timing and received-power consistency; reflections from aircraft, terrain or controlled targets can then test bistatic detection across several carriers and bandwidths. Each carrier must be treated as a separate waveform and propagation case, with transmitter power, antenna sector, scheduling, synchronization and local RFI recorded as metadata.

This converts the cellular infrastructure into a distributed test field for the wideband multiradome radar concept. It does not turn the system into an active radar: the radomes remain passive receivers, and the useful observables come from known transmissions, their geometry and their reflected components.

\begin{figure}[H]
\centering
\includegraphics[width=0.98\textwidth]{en/fig15_aircraft_two_radomes_100km.png}
\caption{Two-node aircraft scenario: the two enlarged radomes define the nominal 100 km baseline, the aircraft is above and offset from it, and the blue/red paths and cones represent independent received signals and angular observations used for power, AOA, TDOA and Doppler consistency.}
\label{fig:fig15_aircraft_two_radomes_100km}
\end{figure}

\begin{figure}[H]
\centering
\includegraphics[width=0.95\textwidth]{en/fig09_fluxo_processamento.png}
\caption{Multistatic passive-processing flow: reference and surveillance inputs converge into channel equalization, direct-path cancellation, delay--Doppler processing, controlled detection and association; lower branches represent emitter priors and the resulting state estimate.}
\label{fig:fig09_fluxo_processamento}
\end{figure}
